% 考古天文学（综述）
% license CCBYNCSA3
% type Wiki

本文根据 CC-BY-SA 协议转载翻译自维基百科\href{https://en.wikipedia.org/wiki/Archaeoastronomy}{相关文章}。

\begin{figure}[ht]
\centering
\includegraphics[width=6cm]{./figures/b9737bc319a7b6fe.png}
\caption{冬至之日，旭日东升，阳光仅在此时照亮爱尔兰纽格兰奇（Newgrange）墓冢的内部石室。} \label{fig_ArmY_1}
\end{figure}
考古天文学（Archaeoastronomy，又作 Archeoastronomy）是一门跨学科（interdisciplinary）\(^\text{[1]}\)或多学科（multidisciplinary）\(^\text{[2]}\)的研究领域，主要探讨古代人类如何“理解天空中的现象、如何利用这些现象，以及天空在其文化中所扮演的角色”\(^\text{[3]}\)。Clive Ruggles 指出，将考古天文学简单地视为“古代天文学的研究”是误导性的，因为现代天文学是一门科学学科，而考古天文学则关注不同文化对天象的象征性与文化性解释\(^\text{[4][5]}\)。考古天文学常与民族天文学（ethnoastronomy）并行研究，后者是一门人类学分支，专门研究当代社会中的观星活动与文化观念。此外，考古天文学还与历史天文学（historical astronomy）密切相关—— 后者利用古代的天象记录来解决天文学问题，以及天文学史（history of astronomy），即通过书面文献来评估古代天文学实践的发展历程\(^\text{[6]}\)。
\begin{figure}[ht]
\centering
\includegraphics[width=6cm]{./figures/5d3ecc0dd8759825.png}
\caption{从西西里岛丰达凯利－凡蒂纳（Fondachelli Fantina）的史前遗址皮佐·文托（Pizzo Vento）观测到的春分或秋分日落景象。} \label{fig_ArmY_2}
\end{figure}
考古天文学（Archaeoastronomy）使用多种方法来揭示古代天文实践的证据，这些方法包括考古学、 人类学、 天文学、 统计学与概率论、以及历史学\(^\text{[7]}\)。由于这些方法来源多样、数据类型差异巨大，如何将它们整合为一致而连贯的论证，一直是考古天文学家面临的长期挑战\(^\text{[8]}\)。考古天文学在景观考古学（landscape archaeology）和认知考古学（cognitive archaeology）中占据互补地位。物质证据及其与天空的关联可以揭示更广阔的自然景观如何融入信仰体系，例如玛雅天文学（Mayan astronomy）与农业周期之间的关系\(^\text{[8]}\)。其他结合认知与景观概念的研究实例包括：探讨定居点道路布局中所体现的宇宙秩序（cosmic order）\(^\text{[10][11]}\)。  

考古天文学可适用于所有文化与各个历史时期。尽管不同文化对天空的意义理解差异显著，但在研究古代信仰时，仍然可以跨文化地应用一套科学方法\(^\text{[12]}\)。或许正是因为必须在社会学维度与科学方法之间取得平衡，Clive Ruggles 才将考古天文学形容为—— “一个在一端产出高质量学术研究，而在另一端则存在近乎疯狂的无拘臆测的领域”\(^\text{[13]}\)。
\subsection{历史}
在他所撰写的《天文考古学（Astro-archaeology）》简史中，约翰·米切尔（John Michell）指出，过去两个世纪以来，对古代天文学的研究地位经历了这样的转变：“从疯狂（lunacy）到异端（heresy），再到有趣的想法（interesting notion），最终到达正统的大门（the gates of orthodoxy）。”  

近二十年之后，我们仍可以追问：  
考古天文学是否依然停留在正统的大门前，  
还是已经迈入其间？  

—— 托德·博斯特威克（Todd Bostwick）引述约翰·米切尔\(^\text{[14]}\)

两百年前，当约翰·米切尔尚未写下上述文字时，世界上尚无“考古天文学家”，也尚未出现“职业考古学家”，但已有天文学家与古物学家（antiquarians）。他们的一些研究工作被视为考古天文学的先驱：例如古物学家们注意到英格兰乡间遗迹的天文朝向，正如威廉·斯图克利（William Stukeley）在1740年对巨石阵（Stonehenge）所做的研究\(^\text{[15]}\)；约翰·奥布里（John Aubrey，1678）\(^\text{[16]}\)与亨利·尚西（Henry Chauncy，1700）\(^\text{[17]}\)则在教堂的朝向中寻找类似的天文学原则。十九世纪晚期，理查德·普罗克特（Richard Proctor）与查尔斯·皮亚齐·史密斯（Charles Piazzi Smyth）等天文学家研究了金字塔的天文朝向\(^\text{[18]}\)。

“考古天文学”（Archaeoastronomy）这一术语由伊丽莎白·切斯利·贝蒂（Elizabeth Chesley Baity）于1973年提出（依循欧安·麦基〔Euan MacKie〕的建议）\(^\text{[19][20]}\)。然而，作为一门研究主题，其历史或许更为久远，这取决于对“考古天文学”概念的定义。  

克莱夫·拉格尔斯（Clive Ruggles）\(^\text{[21]}\)指出，十九世纪中期的海因里希·尼森（Heinrich Nissen）或许可以被视为最早的考古天文学家。罗尔夫·辛克莱尔（Rolf Sinclair）\(^\text{[22]}\)认为，十九世纪末至二十世纪初的诺曼·洛克耶（Norman Lockyer）可被称为“考古天文学之父（father of archaeoastronomy）”。欧安·麦基（Euan MacKie）\(^\text{[23]}\)则将其起源界定得更晚，他写道：“……考古天文学的诞生与现代繁荣，无疑应归功于亚历山大·汤姆（Alexander Thom）在1930年代至1970年代间于英国的研究工作。”
\begin{figure}[ht]
\centering
\includegraphics[width=6cm]{./figures/204a1de1a23caa92.png}
\caption{早期的考古天文学家在不列颠群岛进行巨石遗迹（Megalithic constructs）的测量研究，例如位于伦敦德里郡（County Londonderry）的奥格利什（Auglish）遗址，试图通过统计分析发现其中的规律性模式。} \label{fig_ArmY_3}
\end{figure}
在20世纪60年代，工程师亚历山大·汤姆（Alexander Thom）与天文学家杰拉尔德·霍金斯（Gerald Hawkins）的研究重新激发了人们对古代遗址天文学特征的兴趣。霍金斯曾提出巨石阵（Stonehenge）是“一台新石器时代的计算机”\(^\text{[24]}\)。尽管霍金斯的主张大多被否定\(^\text{[25]}\)，但汤姆的研究却不同凡响——他对不列颠群岛巨石遗址的测量结果，提出了当时广泛存在精确天文学实践的假说\(^\text{[26]}\)。欧安·麦基（Euan MacKie）意识到汤姆的理论需要通过实地检验，于1970年与1971年在阿盖尔郡（Argyllshire）的金特劳（Kintraw）独立石遗址展开发掘，以验证汤姆预测的石碑上方山坡上是否存在观测平台。结果确实发现了一个人工平台——这一发现似乎印证了汤姆关于长线天文对齐（long alignment）的假说（金特劳被认为是一个精确的冬至观测点）。麦基随后又在艾莱岛（Islay）的卡尔图恩石圈（Cultoon stone circle）检验了汤姆的几何理论，同样得到了积极结果。麦基因此大体接受了汤姆的结论，并据此发表了关于不列颠的新史前学著作\(^\text{[27]}\)。然而，克莱夫·拉格尔斯（Clive Ruggles）在重新评估汤姆的实地工作后认为，汤姆关于“高度精确天文学”的主张并未得到充分证据支持\(^\text{[28]}\)。尽管如此，汤姆的影响依然深远。埃德温·C·克鲁普（Edwin C. Krupp）在1979年写道：“几乎凭一己之力，他确立了考古天文学实地研究与解释的标准，而他令人惊叹的成果在过去三十年间引发了广泛争议。”\(^\text{[29]}\)他的影响力至今不衰，对观测数据进行统计检验的做法仍然是考古天文学的重要研究方法之一\(^\text{[30][31]}\)。
\begin{figure}[ht]
\centering
\includegraphics[width=6cm]{./figures/b822574214449a3d.png}
\caption{人们提出过这样的观点：玛雅遗址（如乌斯马尔 Uxmal）可能是按照天文对齐（astronomical alignments）的原则建造的。} \label{fig_ArmY_4}
\end{figure}
在“新大陆”（New World）的研究路径中，人类学家更充分地考虑了天文学在美洲原住民文明（Amerindian civilizations）中的作用，这一方法与旧大陆的研究方式存在显著差异。他们能够利用欧洲史前学所缺乏的资料来源，例如民族志记录\(^\text{[32][33]}\)以及早期殖民者留下的历史文献。在安东尼·阿韦尼（Anthony Aveni）\(^\text{[34][35]}\)的开创性研究启发下，新大陆的考古天文学家能够提出有关动机（motives）的解释，而在旧大陆，这类解释往往只能停留在推测层面。然而，由于新大陆研究侧重于历史资料，其关于天文观测“高精度”的某些论断，在与欧洲统计学导向的研究相比时显得相对薄弱\(^\text{[36]}\)。

这一分歧在1981年国际天文学联合会（IAU）于牛津召开的会议上达到顶点\(^\text{[37]}\)。会议参与者的方法论和研究问题差异极大，以至于会议论文被分为两卷分别出版\(^\text{[38][39]}\)。尽管如此，会议被认为成功地促进了研究者之间的交流与合作，自那以后，牛津系列会议每隔四至五年便会在世界各地持续举办。随后的会议推动了考古天文学向更加跨学科的方向发展。研究者们开始尝试将考古学的情境性（contextuality）与天文学结合\(^\text{[40]}\)——这大体描述了当代考古天文学的现状。如今的考古天文学不再仅仅致力于“证明古代天文学的存在”，而是进一步追问：人们为何会对夜空产生兴趣？
\subsection{与其他学科的关系}
“……考古天文学最令人喜爱的特质之一，  
就是它能让来自不同学科的学者彼此争论不休。” 
 
—— 克莱夫·拉格尔斯（Clive Ruggles）\(^\text{[41]}\)

考古天文学长期以来被视为一门跨学科（interdisciplinary）的研究领域，它结合书面与非书面证据，以探究其他文化的天文学。因此，它可以被视为连接多种研究古代天文学路径的桥梁：包括天文考古学（astroarchaeology，现已废弃的术语，指通过古代建筑或地景的朝向推导天文信息的研究）、天文学史（history of astronomy，主要研究书面文献证据），以及民族天文学（ethnoastronomy，借助民族历史文献与当代民族志研究）\(^\text{[42][43]}\)。

考古天文学作为一门跨学科领域的发展也体现在研究者的背景多样性上。该领域的研究由受训于不同学科的学者共同开展。近期多篇博士论文的作者将其研究描述为涉及考古学与文化人类学；也包括历史学的多个分支——如特定地区与时期的历史、科学史与宗教史；同时探讨天文学与艺术、文学、宗教之间的关系。只有极少数研究者将其工作直接归类为“天文学”，即便如此，也通常仅作为次要范畴\(^\text{[44]}\)。

无论是实践中的考古天文学家，还是观察该领域的学者，他们对考古天文学的理解方式各不相同。部分研究者将其与科学史联系起来，关注某一文化如何观察自然，并建立概念体系以为这些观察赋予秩序\(^\text{[45]}\)；另一些则关注推动特定历史人物采用某种天文学概念或技术的政治动机\(^\text{[46][47]}\)。艺术史学者理查德·波斯（Richard Poss）采取了更为灵活的立场，认为北美西南地区的天文岩画（astronomical rock art）应当依据“西方艺术史与艺术批评的诠释学传统”进行解读\(^\text{[48]}\)。然而，天文学家关注的问题有所不同 —— 他们希望为学生揭示本学科的可识别先驱，并尤为关心一个核心问题：如何确认某个遗址确实是出于天文目的而有意设计的\(^\text{[49]}\)。

考古学界的反应（The reactions of professional archaeologists to archaeoastronomy）专业考古学界对考古天文学（archaeoastronomy）的反应始终褒贬不一。一些学者表现出难以理解甚至敌意的态度——这种敌意的形式多样，既包括主流考古学界将其视为“边缘学科”而拒绝接纳，也包括由于考古学家注重文化背景而早期考古天文学家偏重定量方法，导致双方沟通困难、理念分歧\(^\text{[50]}\)。然而，近年来考古学界逐渐吸收了许多来自考古天文学的研究成果与洞见，并将其纳入考古学教材\(^\text{[51]}\)。如前所述，一些学生甚至以考古天文学为主题撰写博士论文，显示出该领域在学术界的影响力正逐步扩大。

由于考古天文学家对这一学科的性质与边界意见分歧极大，他们甚至在“该如何命名”这一问题上存在争论。三大国际学术组织均将考古天文学视为研究文化与天文学关系的学科，采用“文化中的天文学”（Astronomy in Culture）或其译名作为正式术语。迈克尔·霍斯金（Michael Hoskin）认为，该领域的重要部分在于“事实收集”而非“理论建构”，因此建议将这一方向命名为“考古地形学”（Archaeotopography）\(^\text{[52]}\)。克莱夫·拉格尔斯（Clive Ruggles）与尼古拉斯·桑德斯（Nicholas Saunders）则提出，应使用“文化天文学”（Cultural Astronomy）作为总称，以统一不同民俗天文学研究方法\(^\text{[53]}\)。另一些学者认为，“天文学”一词并不准确，他们主张研究的核心是“宇宙观”（cosmology）而非“天文学”（astronomy），并建议采用西班牙语术语“cosmovisión”以避免“logos”的科学理性含义\(^\text{[54]}\)。

当研究争论围绕方法论极端化时，人们常用一种“颜色代号”来区分不同学派 —— 这一命名方式源自第一次牛津会议（Oxford Conference）论文集的装帧颜色，该会议首次明确区分了不同的研究取向\(^\text{[55]}\)。“绿色学派”（Green，旧大陆）考古天文学家主要依赖统计方法，常被批评忽视天文行为的文化与社会语境；而“棕色学派”（Brown，新大陆）学者则拥有丰富的民族志与历史资料，但在测量与统计分析方面常被指“过于随意”\(^\text{[56]}\)。自20世纪90年代初以来，如何在这两种方法间寻求整合与平衡，一直是考古天文学界的重要讨论议题\(^\text{[57][58]}\)。
\subsection{方法论}
“长期以来，我一直相信，这样的多样性必然需要一种包罗万象的理论来解释。  
但我现在认为，当初那样的想法实在太天真了。”

—— 斯坦尼斯瓦夫·伊瓦尼舍夫斯基（Stanislaw Iwaniszewski）\(^\text{[59]}\)

考古天文学（archaeoastronomy）并没有唯一的研究方法。学者之间的分歧往往并非出现在“自然科学家”与“社会科学家”之间，而是取决于研究者所能获得的数据类型或研究地区的不同。在“旧大陆”（Old World），研究者往往仅能依靠遗址本身获取信息；而在“新大陆”（New World），则可结合民族志（ethnography）与历史记录来补充。这种地域性分化使得考古天文学在不同地区的发展路径独立而多样，其影响至今仍可在各地研究中观察到。总体而言，研究方法可大致分为两类，不过许多近年的项目已开始将两种方法结合使用，以实现更全面的解释框架。
\subsubsection{绿色考古天文学}
“绿色考古天文学”一名来源于《旧大陆考古天文学》（Archaeoastronomy in the Old World）一书的绿色封面\(^\text{[60]}\)。这种方法主要以统计学为基础，特别适用于史前遗址研究—— 在这些遗址中，社会文化证据相较于有文字记载的历史时期较为稀少。该方法的基本技术由亚历山大·汤姆（Alexander Thom）在其对英国巨石遗址的广泛测量中建立。

汤姆的核心研究目标是检验史前人群是否掌握了高精度的天文学。他认为，通过“地平线天文学”（horizon astronomy），观察者能够以日为单位，精确估算一年中的特定日期。这种观测方法依赖寻找一个特定地点，在某一天太阳会刚好落入地平线的某个凹口或山间缺口中。常见的例子是：太阳被山峰遮挡，但在某个特定日期，太阳会在山的另一侧重新露出极小一部分，形成所谓的“双重日落”（double sunset）现象。下方的示意动画展示了一个假想遗址的两次日落情景：一次发生在夏至前一天，另一次则正好是夏至当天。在后者中，太阳出现了典型的“双重日落”现象。

为了验证这一设想，汤姆（Thom）对数百处石阵（stone rows）与石环（stone circles）进行了测量。单个排列的方向可能纯属偶然，但他计划通过统计显示，这些排列的整体分布并非随机，从而证明至少有一部分排列的朝向具有明确的天文意图。他的研究结果显示，古人或许将一年划分为八、十六，甚至三十二个大致相等的时间区段。其中包括两个至日（solstices）、两个分点（equinoxes），以及四个“交叉季度日”（cross-quarter days）—— 即位于至日与分点之间的中点日。这些节气日期与中世纪凯尔特历法（Celtic calendar）中的节日相对应。\(^\text{[61]}\)虽然并非所有这些结论都被后世学者完全接受，但汤姆的研究对考古天文学（archaeoastronomy）产生了深远而持久的影响，尤其是在欧洲地区。

尤安·麦基（Euan MacKie）支持汤姆的分析，他通过将新石器时代的英国与玛雅文明进行比较，为汤姆的研究加入了考古学背景，提出当时的社会可能是分层结构的。\(^\text{[27]}\)为了验证这些假设，麦基在苏格兰的几处史前观测台遗址进行了发掘。例如金特劳（Kintraw）遗址以一块四米高的竖石（standing stone）而闻名。汤姆认为，这块石头正好指向远处地平线上位于尤拉岛（Jura）上两座山 —— 贝恩·夏尼德（Beinn Shianaidh）与贝恩·欧·凯利亚斯（Beinn o’Chaolias）——之间的缺口。\(^\text{[62]}\)他推测，在冬至时，太阳会在这一地平线缺口处出现“双重日落”（double sunset）现象。然而，从地面高度观察，这一景象会被地形中的山脊遮挡，观测者必须抬高约两米的高度——因此需要另一个观测平台。麦基在峡谷对岸确实发现了一个由小石块堆成的平台。虽然缺乏文物遗迹令一些考古学家质疑其用途，而岩相分析（petrofabric analysis）也未能得出决定性结论，但随后在梅斯豪（Maes Howe）\(^\text{[63]}\)与布什丘陵金菱形饰板（Bush Barrow Lozenge）\(^\text{[64]}\)上的进一步研究让麦基得出结论：虽然用“科学”（science）一词来形容史前观测或许过于时代错置，但汤姆关于“高精度对齐”的整体观点基本正确。\(^\text{[65]}\)

与此相对，克莱夫·拉格尔斯（Clive Ruggles）则认为汤姆的调查在数据选择上存在问题。\(^\text{[66][67]}\)也有学者指出，“地平线天文学”的精度受到近地平线折射变化的限制。\(^\text{[68]}\)对“绿色考古天文学”（Green archaeoastronomy）更深层的批评在于：虽然它能回答“古人是否对天文学感兴趣”，但由于缺乏社会文化层面的分析，它难以解释“古人为何对此感兴趣”。正如人类学家基思·金蒂格（Keith Kintigh）所言：“直白地说，在许多情况下，某一考古天文学主张是否正确，对人类学的发展并无太大意义，因为这些信息无法回应当下的解释性问题。”\(^\text{[69]}\)尽管如此，“对齐研究”（study of alignments）仍然是考古天文学的重要组成部分，尤其在欧洲地区，依然是该领域研究的核心内容之一。\(^\text{[70]}\)
\subsubsection{棕色考古天文学}
与以统计分析和天文对齐（alignment）为主的“绿色考古天文学”相对，“棕色考古天文学”被认为更接近于天文学史（history of astronomy）或文化史（cultural history）。它依赖历史文献与民族志（ethnographic）资料，以深化对早期天文学及其与历法（calendars）与宗教仪式（ritual）的关系之理解。\(^\text{[55]}\)由于西班牙殖民编年史家与民族学研究者留下了大量关于本土习俗与信仰的记录，“棕色考古天文学”往往与美洲古代天文学的研究联系在一起。\(^\text{[71][72][32][33]}\)  

最著名的一个案例是奇琴伊察（Chichen Itza）遗址。与其通过分析遗址的结构来寻找天文对齐的方向，棕色考古天文学家更倾向于先从民族志资料中考察 —— 玛雅人（Mayans）认为哪些天象具有重要意义，再在考古遗迹中寻找与之对应的物质证据。一个若无历史记录可能被忽视的重要例子是玛雅人对金星（Venus）的兴趣。这种兴趣由《德累斯顿手抄本》（Dresden Codex）所证实，其中包含金星在天空中出现与消失的详细表格。\(^\text{[73]}\)这些周期被赋予占星与宗教仪式的意义，因为金星被视为羽蛇神奎兹尔科亚特尔（Quetzalcoatl）或肖洛特尔（Xolotl）的象征。\(^\text{[74]}\)在奇琴伊察（Chichen Itza）、乌斯马尔（Uxmal），以及其他一些中美洲（Mesoamerican）遗址中，都能发现建筑特征与金星升落位置的对应关系。\(^\text{[75]}\)
\begin{figure}[ht]
\centering
\includegraphics[width=8cm]{./figures/9fcd5f4746017402.png}
\caption{卡拉科尔”（El Caracol），奇琴伊察（Chichen Itza）的一座可能用于天文观测的神庙。} \label{fig_ArmY_5}
\end{figure}
战士神庙（The Temple of the Warriors）上有描绘羽蛇（feathered serpents）的图像，这些形象与羽蛇神——奎兹尔科亚特尔（Quetzalcoatl）或库库尔坎（Kukulcan）——有关。因此，这座建筑朝向地平线上金星（Venus）傍晚初现之处（此时恰逢雨季）的定向，可能具有特殊的象征意义。\(^\text{[76]}\)然而，由于这一事件的日期与方位角（azimuth）会持续变化，这一朝向更可能与太阳有关，而非金星。\(^\text{[77]}\)

安东尼·阿维尼（Anthony Aveni）指出，奇琴伊察（Chichen Itza）中另一座与金星（以库库尔坎形象）及雨季相关的建筑是卡拉科尔塔（El Caracol）。\(^\text{[78]}\)这是一座带有圆形塔楼的建筑，其门朝向四个基本方向。塔基朝向金星在北方最偏西落下的位置。此外，上层平台的柱基（stylobate）被涂成黑色与红色，这两种颜色象征着金星作为晨星与昏星的双重身份。\(^\text{[79]}\)然而，塔楼的窗户不过是狭缝结构，几乎不能有效采光，但十分适合用于向外观测。\(^\text{[80]}\)在探讨考古天文学遗址可信度时，科特（Cotte）与拉格尔斯（Ruggles）指出，将卡拉科尔解释为“天文观测台”的观点在专家间仍存争议，其可信度大约处于他们划分的四个等级中的第二级。\(^\text{[81]}\)

阿维尼认为，“棕色方法论”（Brown methodology）的优势之一在于，它能够揭示那些统计分析无法捕捉的天文学体系。他以印加帝国（Inca Empire）的天文学为例。印加人以库斯科（Cusco）为中心，用“赛克线”（ceques）辐射划分整个帝国的空间结构。这些线在各个方向上形成大量排列，乍看之下似乎缺乏天文意义。然而，民族史料表明，这些方向确实具有宇宙学（cosmological）与天文学（astronomical）象征，其上的各个地景点在一年中不同时间具有不同的宗教或历法重要性。\(^\text{[82][83]}\)在东亚地区，考古天文学的发展源自天文学史（history of astronomy），许多研究旨在寻找文献记载的天象之物质对应物。这是因为东亚，特别是中国，拥有丰富的天文历史记录，其连续性可追溯至公元前二世纪的汉代。\(^\text{[84]}\)

这一方法的批评在于其统计基础较弱。沙弗尔（Schaefer）尤其质疑卡拉科尔建筑中所声称的天文对齐的可靠性。\(^\text{[85][86]}\)由于考古天文学涉及极为广泛的证据类型——包括遗址与器物—— 因此并不存在单一统一的研究路径。\(^\text{[87]}\)尽管如此，学界普遍认为考古天文学并非孤立学科。作为一门跨学科（interdisciplinary）的研究领域，其研究结论必须在考古学与天文学两个层面上均具有解释力。若研究采用考古学中常见的理论工具，如类比（analogy）与同源（homology），并且能体现出对天文学中“精度与准确性”的理解，其成果更容易被视为可靠。近年的系统性研究已将定量分析（quantitative analyses）与民族志类比（ethnographic analogies）及其他情境性证据结合，应用于玛雅地区（Maya area）\(^\text{[88]}\)以及中美洲其他地区（Mesoamerica）\(^\text{[89]}\)的建筑定向研究。
\subsection{资料来源}
由于考古天文学（Archaeoastronomy）研究的是人类与天空互动的多种方式，  
因此其信息来源极为多样，涉及多种与天文实践相关的证据形式。
\subsubsection{方位}
考古天文学中最常见的数据来源之一是对遗址“方位对齐”（alignment）的研究。这一研究基于这样的假设：考古遗址的轴线方向（axis of alignment）具有天文学上的意义，即指向特定的天体目标（astronomical target）。  

“棕色考古天文学家”（Brown archaeoastronomers）通常通过阅读历史或民族志文献（historical / ethnographic sources）来验证这一假设；而“绿色考古天文学家”（Green archaeoastronomers）则倾向于通过统计方法证明方位不太可能是随机选择的，例如展示多个遗址中存在共同的定向模式（common patterns of alignment）。“方位”通常通过测量结构物的方位角（azimuth）与其面对地平线的高度角（altitude）而确定。\(^\text{[90]}\)方位角是指该结构相对于正北方向的夹角。通常使用经纬仪（theodolite）或罗盘（compass）来测量方位角。罗盘使用方便，但必须考虑地球磁场偏离真北的“磁偏角”（magnetic declination）。此外，在存在磁干扰（magnetic interference）的环境中——例如有金属脚手架支撑的遗址 —— 罗盘的可靠性较差。而且罗盘的测量精度仅约为半度（0.5°）。\(^\text{[91]}\)

若正确使用，经纬仪的精度远高于罗盘，但同时也更难操作。经纬仪本身无法直接与真北对齐，因此必须通过天文观测来校准刻度，通常使用太阳的位置（the position of the Sun）进行校准。\(^\text{[92]}\)由于天体的位置会随地球自转而改变，这些校准观测的时间必须被精确记录，否则测量结果会产生系统性误差（systematic error）。地平线高度角（horizon altitude）可使用经纬仪或倾角仪（clinometer）来测量。
\subsubsection{文物}
\begin{figure}[ht]
\centering
\includegraphics[width=6cm]{./figures/a6185b053b471298.png}
\caption{安提凯希拉机械装置（主残片）} \label{fig_ArmY_6}
\end{figure}
以“内布拉天盘”（Sky Disc of Nebra）为例——据称是一件描绘宇宙的青铜时代文物\(^\text{[93][94]}\)—— 其分析方法与考古学中其他分支的常规出土后分析（post-excavation analysis）类似。研究者会对文物进行详细检查，并尝试与其他民族的历史或民族志记录进行类比。找到的平行案例越多，该解释被其他考古学家接受的可能性就越高。

一个更日常的例子来自罗马帝国时期出土的鞋履与凉鞋，其上刻有占星符号（astrological symbols）。尽管鞋履的使用功能众所周知，但考古学家 Carol van Driel-Murray 提出，这些刻印符号可能赋予鞋履宗教或医疗意义（spiritual or medicinal meanings）。\(^\text{[95]}\)这一观点得到了其他史料的支持——当时占星符号与医疗实践的密切联系以及相应的历史记载都为这种解释提供了依据。\(^\text{[96]}\)

另一件著名的具有天文用途的文物是安提凯希拉机械装置（Antikythera mechanism）。
通过对该装置的结构分析，并参考西塞罗（Cicero）对类似仪器的描述，可以合理推测其用途为一种用于计算天体位置的仪器。这一推测还得到了装置上符号标识的进一步支持，这些符号使得机械盘的读数成为可能。\(^\text{[97]}\)
\subsubsection{艺术与铭文}
\begin{figure}[ht]
\centering
\includegraphics[width=6cm]{./figures/ac053adce81cdd2d.png}
\caption{法哈达丘岩刻上“太阳匕首”（Sun Dagger）在不同日期的位置示意图} \label{fig_ArmY_7}
\end{figure}
艺术与铭文不仅限于独立文物，也可能出现在考古遗址的墙壁或岩面上，以绘画或雕刻的形式存在。有时铭文的内容足够直接，甚至能说明遗址的用途。例如，在伊塔诺斯（Itanos）发现的一块希腊石碑（stele）上，铭文被译为：“庇护者为宙斯·埃波普西奥斯（Zeus Epopsios）树立此碑。冬至。若有人想知道：从‘小猪’与此碑起，太阳转向。”\(^\text{[98]}\)在中美洲，玛雅与阿兹特克文明留下了玛雅与阿兹特克抄本（codices）。这些是用加工树皮纸 Amatl 制成的折叠书册，上面以玛雅或阿兹特克文字刻绘图符。其中最著名的《德累斯顿抄本（Dresden Codex）》包含关于金星周期的资料，再次确认了金星在玛雅文化中的重要地位。\(^\text{[73]}\)

更具挑战性的是那些阳光与阴影在不同季节与时间与岩刻互动的情况。\(^\text{[99]}\)一个广为人知的例子是法哈达丘（Fajada Butte）上的“太阳匕首”（Sun Dagger），在此处，一道阳光穿过岩缝，正好扫过一枚螺旋形岩刻。\(^\text{[100]}\)阳光在螺旋上的位置随季节变化：在夏至时，一道“光匕首”穿过螺旋中心；在冬至时，两道光线出现在螺旋两侧。研究者推测，这一岩刻可能是为了标示这些天文事件而制作的。最近的研究发现，在美国西南部与墨西哥西北部存在许多类似遗址。\(^\text{[101][102]}\)学者们指出，这些遗址中太阳在至日（solstice）时的光影标记数量之多，在统计上表明这些设计极可能是有意用于标记冬夏至。\(^\text{[103]}\)位于新墨西哥州查科峡谷（Chaco Canyon）的法哈达丘“太阳匕首”遗址尤其突出，因为它的光影标记精确记录了太阳与月球周期中的关键节点：夏至、冬至、春秋分，以及月球 18.6 年大与小停点（major/minor lunar standstills）。\(^\text{[104][105]}\)此外，在法哈达丘的另外两个地点，还发现了五处光影标记，分别对应夏至、冬至、春秋分与太阳正午。\(^\text{[106]}\)查科峡谷及其周边地区的许多大型建筑与建筑群之间的排列方向，也与“太阳匕首”遗址上所标示的太阳与月亮方位保持一致。\(^\text{[107]}\)

然而，如果没有任何民族志或历史资料支持这种解释，其接受度就取决于——北美是否存在足够多的岩刻遗址，以至于这种“光影对应”可以单纯地归因于巧合。当岩刻仍与现存民族存在文化延续时，民族天文学家（ethnoastronomers）便可以直接向当代传承者询问符号的意义，以帮助揭示这些天文图像背后的真实文化意图。
\subsubsection{民族志}
除了古代人自己留下的遗迹之外，还有那些与他们接触过的外来者的记录。例如，西班牙征服者（Conquistadores）的历史文献就是了解前哥伦布时期美洲人的重要资料来源。此外，民族志学者（ethnographers）也提供了关于世界上许多其他民族的丰富记录。

天文学家安东尼·阿维尼（Anthony Aveni）以“天顶日（zenith passages）”的重要性为例，说明民族志在考古天文学中的作用。对于生活在北回归线与南回归线之间的民族而言，
每年有两天正午时分太阳会正好经过头顶（天顶），此时地面上不会产生影子。在中美洲的某些地区，这一天被视为极其重要，因为它预示着雨季的到来，与农业周期密切相关。直到今天，生活在中美洲的玛雅印第安人依然认为这一知识具有重要意义。民族志资料启发考古天文学家推测：这一天在古代玛雅社会中也可能被赋予特殊意义。在像蒙特阿尔班（Monte Albán）和索奇卡尔科（Xochicalco）这样的遗址中，发现了被称为“天顶管（zenith tubes）”的竖井，当太阳经过头顶时，阳光正好照亮地下室。只有通过民族志记录，我们才能推测：在玛雅社会中，这种阳光照射的时刻被视为具有重要的仪式或历法意义。\(^\text{[108]}\)此外，有学者声称，一些遗址的建筑朝向与太阳在天顶日出、日落时的方位一致。然而，研究表明：能够与这些天象相对应的方位非常有限，因此它们更可能有其他解释。\(^\text{[109]}\)

民族志同时也对过度阐释遗址提出警示。例如，在查科峡谷（Chaco Canyon）的一处遗址，
发现了一个描绘“星星、弯月与手印”的岩画。一些天文学家认为，这幅图记录了公元1054年的超新星爆发。\(^\text{[110]}\)然而，对类似“超新星岩刻（supernova petroglyph）”的重新研究对这种解释提出了质疑。\(^\text{[111]}\)科特（Cotte）与拉格尔斯（Ruggles）在其研究中将这幅“超新星岩刻”列为已被完全驳斥的案例\(^\text{[81]}\)，而人类学证据则表明了其他文化性解释。例如，与查科遗址有强烈祖源联系的祖尼族（Zuni），在他们的观日台上刻有“弯月、星星、手印与太阳圆盘”的符号，与查科的岩画非常相似。\(^\text{[112]}\)

民族天文学（Ethnoastronomy）不仅在美洲重要，在其他地区也发挥着关键作用。例如，关于澳大利亚原住民的民族学研究，揭示了大量他们的本土天文学知识（Indigenous astronomies），\(^\text{[113][114]}\)以及他们与现代社会在天文认知与世界观上的互动。\(^\text{[115]}\)
\subsection{重现古代的星空}
“……尽管不同文化中确实出现了不同的科学方法与结果，但这并不能支持那些质疑科学是否有能力对我们所处世界作出可靠陈述的观点。”

—— 斯蒂芬·麦克拉斯基（Stephen McCluskey）\(^\text{[116]}\)

当研究者获得了可供验证的数据后，通常需要尝试重建古代的天象条件，以便将这些数据置于其历史语境之中进行分析。
\subsubsection{赤纬}
为了计算一处建筑或遗迹面向的天体方位，必须使用一种坐标系统，而恒星正好提供了这样的系统。在晴朗的夜晚，人们可以观察到恒星围绕天极（celestial pole）旋转。这个旋转中心即是：北天极（North Celestial Pole），位于 +90°；南天极（South Celestial Pole），位于 −90°。\(^\text{[117]}\)恒星在天空中划出的同心圆，即是天球纬线（celestial latitude），称为赤纬线（declination lines）。连接东、西方地平点（在地平面平坦的情况下）并且位于南北天极之间中点的那条弧线，称为天球赤道（Celestial Equator），其赤纬为 0°。可见的赤纬范围取决于观测者所在的地球位置：
只有位于地球北极的观察者，在夜晚完全看不到南天球的恒星（如下图所示）。一旦确定了建筑面向地平线上某一点的赤纬值，研究者便可以判断 —— 在这一方向上，是否能看到特定的天体（例如太阳、月亮或某颗恒星）。
\begin{figure}[ht]
\centering
\includegraphics[width=10cm]{./figures/23cec6286466a492.png}
\caption{不同纬度下可见天空区域示意图} \label{fig_ArmY_8}
\end{figure}
\subsubsection{太阳定位}
恒星的赤纬位置是固定的，而太阳的赤纬却不是。太阳的升起点在一年中会不断变化，
像钟摆一样在两个极限点之间来回摆动 —— 这两个极限点分别由夏至与冬至（solstices）标记。当太阳到达极限点时移动最慢，而在中点（春分与秋分附近）时则移动得最快。因此，如果一位考古天文学家根据方位角（azimuth）与地平高度（horizon height）计算出某遗址的观测方向对应的赤纬为 +23.5°，那么他/她就可以推断该遗址是为了观测夏至日出而建，而无需真正等到 6 月 21 日 去验证这一点。\(^\text{[118]}\)（详见主条目：太阳观测史 History of Solar Observation）
\subsubsection{月亮定位}
月球的运行要比太阳复杂得多。与太阳类似，月球的视运动也在两个极限之间往复，这两个极限称为月极（lunistices），相当于太阳的“至点（solstices）”。但不同之处在于：月球在这两个极限之间完成一次往返只需一个恒星月（sidereal month），而太阳需要整整一年。更复杂的是，月球运动极限（即月极）本身也在18.6 年周期内变化。在大约 9 年多的时间里，月球升落的极限点超出太阳的日出范围；而在接下来的半个周期中，月球的升落点则始终不超过太阳日出的极限范围。此外，古代的月球观测往往更关注月相（phase of the Moon）从一个新月到下一个新月的周期是朔望月（synodic month），这一周期与恒星月完全不同。\(^\text{[119]}\)因此，当研究考古遗址是否具有“月亮相关意义”时，由于月球位置和相位变化极为复杂，相关数据往往显得稀疏或模糊。（详见主条目：月球 Moon）
\subsubsection{恒星定位}
\begin{figure}[ht]
\centering
\includegraphics[width=6cm]{./figures/e54567957947d924.png}
\caption{岁差运动} \label{fig_ArmY_9}
\end{figure}
最后，研究者往往还需要校正恒星的视运动。在整个人类文明的时间尺度上，恒星之间的相对位置几乎保持不变。每个夜晚，恒星似乎围绕着天极（celestial poles）旋转，这源于地球自转的结果。然而，地球的旋转并非完美稳定的。地球像一个陀螺（spinning top）那样，不仅自转，还会轻微晃动（进动，precession）。地球自转轴完成一次完整的晃动大约需要 25,800 年。\(^\text{[120]}\)这种现象对考古天文学家的影响在于：古代恒星的升起位置与今天并不相同，恒星也并非一直围绕北极星（Polaris）旋转。

早在六千多年前，苏美尔人（Sumerians）就注意到了这种轴向运动。他们发现，在其最古老、最南端的城市埃里杜（Eridu），恒星老人星（Canopus）首次能够在南方子午线上空正中天（culminate）。在接下来的几十年间，老人星仍然无法在东北方的邻城乌尔（Ur）被看到，因此，苏美尔人称其为“埃里杜之星（Star of the City of Eridu）”。\(^\text{[121][122]}\)

以埃及金字塔（Egyptian pyramids）为例，研究表明它们曾被精确对准位于天龙座（Draco）中的暗星右枢（Thuban）。\(^\text{[123]}\)从历史尺度看，这种恒星位置的偏移在相对短的时间内就会带来显著变化。例如：在古罗马时期，一个出生于 12 月 25 日的人，其生日那天太阳位于摩羯座（Capricorn）；而在现代，同一天出生的人，太阳则位于射手座（Sagittarius）—— 这正是由于岁差（precession of the equinoxes）造成的结果。
\subsubsection{瞬变天象}
\begin{figure}[ht]
\centering
\includegraphics[width=6cm]{./figures/c6f705e973af98d5.png}
\caption{《巴约挂毯》（Bayeux Tapestry）中描绘的哈雷彗星（Halley’s Comet）} \label{fig_ArmY_10}
\end{figure}
此外，天空中还存在一些瞬变现象，即并非以年度周期发生的天文事件。最容易预测的瞬变事件是日食与月食。就日食（solar eclipses）而言，它们常被用于推算历史事件的确切日期。例如，希罗多德（Herodotus）提到过的一次日食，使我们能够将米堤亚人（Medes）与吕底亚人（Lydians）之间的一场战役（由于这次日食的出现而中止）精确地定年为公元前 585 年 5 月 28 日。\(^\text{[124]}\)

彗星（Comets）中也存在某种规律性，其中最著名的是哈雷彗星（Halley’s Comet）。
然而，作为一类天体，彗星整体上仍具有高度的不可预测性 —— 它们可能在任何时候出现。
一些彗星拥有极长的轨道周期，这使得它们的过去出现和未来回归无法准确推算；另一些彗星可能仅仅掠过太阳系一次，因此从根本上就无法预测。\(^\text{[125]}\)

流星雨（Meteor showers）理论上应当是可预测的，但由于其中许多流星源自彗星碎屑（cometary debris），而这些碎屑的轨道目前尚无法完全计算，因此其出现时间和强度也难以准确预报。\(^\text{[126]}\)古人还曾记录过其他短暂天象，如极光（aurorae）、幻日（sun dogs）以及彩虹（rainbows）等。这些现象与古代天气一样无法预测，但它们在当时的文化或宗教体系中往往被视为重要的天象征兆。
\subsection{考古天文学研究的主要主题}
“天文学在历史上为不同文化群体的生活带来了什么？
答案是多样而丰富的……”

——冯·德尔·钱伯林（Von Del Chamberlain）与 M·简·杨（M. Jane Young）\(^\text{[127]}\)
\subsubsection{天文学与历法的使用}
天文学的重要功能之一，是为了农业生产而制定精确的历法（calendar）。古代文献如**赫西俄德（Hesiod）的《工作与时日》（Works and Days），这本古希腊农事手册部分印证了这一点：在书中，作者将天象观测与自然生态信号（如候鸟迁徙）相结合，用以确定季节的更替。民族天文学（ethnoastronomy）的研究表明，美国西南部的霍皮族（Hopi）会仔细观测太阳的升起与落下位置，以判断何时是播种的适宜时机。\(^\text{[128]}\)然而，对埃塞俄比亚的穆尔西族（Mursi）的研究显示，他们的阴阳合历（luni-solar calendar）相当随意，说明在某些社会中，天文学历法的精度和重要性是有限的。\(^\text{[129]}\)尽管如此，历法仍然几乎是普遍存在于人类社会中的制度，因为它为公共活动的组织与时间调控提供了必要的框架。

在中美洲前哥伦布时期文明中，存在一种名为 Tzolk’in 的260 天历法。它与365 天太阳年历结合使用，形成了一个复杂的历法体系，该体系整合了天文观测与宗教仪式周期。\(^\text{[130]}\)考古天文学的研究表明，整个美索美洲地区的大多数建筑结构都与太阳方位相关，并且结合260 天周期用于确定农耕与仪式活动的时间。在墨西哥南部的湾岸地区（southern Gulf Coast）发现的纪念性宗教建筑群中，通过其方位与时间间隔的分布，可追溯出约公元前 1100 至 700 年期间最早使用这一周期的证据。\(^\text{[131]}\)

古希腊历法形式多样，通常为阴历（lunar calendar），以新月（New Moon）为一个月的起始。然而在实践中，这种历法可能跳日或暂停，导致居民在记日期时会同时标注城邦公历（civic calendar）与“众神之历”（ton theoi）以对照。\(^\text{[132]}\)古希腊缺乏统一历法，这意味着举办泛希腊性活动（如竞技会、宗教仪式）的时间协调十分困难。因此，有学者认为，天文象征可能被用作一种政治上中立的时间参照方式。\(^\text{[133]}\)此外，研究还发现，希腊神庙与拜占庭教堂的朝向测量常常与神祇的命名日（Name Day）、宗教节庆及特殊天象事件有关。\(^\text{[134][135][136]}\)
\subsubsection{神话与宇宙观}
\begin{figure}[ht]
\centering
\includegraphics[width=6cm]{./figures/618fff29adc98147.png}
\caption{约翰内斯·赫维留斯（Johannes Hevelius）于1690年绘制的星座〈船底座（Argo Navis）〉图} \label{fig_ArmY_11}
\end{figure}
研究天空的另一种动机是理解与解释宇宙。在许多古代文化中，神话（myth）是一种实现这一目的的工具。虽然这些解释不符合现代科学的标准，但它们构成了各自的宇宙观。

印加人通过他们帝国的布局来体现其宇宙观。首都库斯科（Cusco）位于帝国的中心，从中心向外放射出称为“塞克线”（ceques）的概念性直线。\(^\text{[137]}\)这些线将帝国中心与四个“苏尤”（suyus）相连，每个苏尤都按照与库斯科的方位（东、南、西、北）来定义。这种“四分宇宙”（quartered cosmos）的观念在整个安第斯地区都普遍存在。人类学家加里·乌尔顿（Gary Urton）在秘鲁安第斯山区的米斯米奈（Misminay）村进行田野调查时发现，这种“四方宇宙”的思想与夜空中银河（Milky Way）的形状密切相关。\(^\text{[138]}\)在一年中的某个季节，银河会斜贯夜空；而在另一个季节，它则会以近乎垂直的方式横切天空。

中国的宇宙布局与“天人秩序”在世界的另一端，对宇宙秩序的观测同样发挥了重要作用。中国北京的紫禁城（The Forbidden City）在布局上体现了古代中国的宇宙秩序（cosmic order）。与安第斯的“四方观”不同，中国古代的宇宙体系包括五个方位：东、西、南、北与中央。\(^\text{[139]}\)紫禁城正位于古代北京的“中宫”位置。从南面朝北进入皇宫，象征着皇帝坐镇北辰之下，众星拱之。这意味着整个天穹都在皇帝的“中心”周围旋转，从而使皇权与宇宙秩序相呼应。这种思想后来通过风水（feng shui）传播到世界各地，成为中国宇宙观的代表形式之一。

关于宇宙运作的观念，也被保存在星座神话（constellation myths）之中。例如，亚马孙的巴拉萨纳族（Barasana）会根据星象来安排他们的年度生活节奏。当他们的星座“毛毛虫-美洲虎（Caterpillar–Jaguar）”（大致对应现代的天蝎座 Scorpius）沉落时，他们便准备捕捉从树上坠落的蛹化毛虫。\(^\text{[140]}\)这些毛虫在食物匮乏的季节中，成为重要的蛋白来源。\(^\text{[141]}\)

阿拉斯加与加拿大北部的“旅行者星座”（The Traveller Constellations）阿拉斯加及加拿大西北部的北德内人（Northern Dene peoples）识别出一组庞大的“旅行者星座”（Traveller constellations），这些星座覆盖了夜空的大部分区域。它们代表着一位文化英雄“旅行者”（The Traveller），根据口头传统，他曾环游世界，后来其灵魂化作夜空中的星辰。\(^\text{[142]}\)这些星座由多个小星群组成，名称采用人体部位的词汇来命名，而银河被认为是“旅行者的足迹”（the Traveller’s trail）。这种星空知识通常通过个人领悟（personal discovery）获得，而非正式教学，并在导航、计时与宗教仪式中扮演重要角色。\(^\text{[143]}\)

希腊与罗马的星座神话（Greek and Roman Constellation Myths）更为人熟知的星座神话来自古希腊与古罗马的文献。这些星座的起源至今仍是学界激烈争论的话题。\(^\text{[144][145]}\)例如，希腊神话中关于“昴宿星团（Pleiades）”中墨洛珀（Merope）消失的故事，可能反映了一次真实的天文事件 —— 某颗昴星的亮度衰退至肉眼不可见。\(^\text{[146]}\)

《哈姆雷特的磨坊》与宇宙神话的学说（Santillana & von Dechend’s Theory）麻省理工学院科学史教授乔治·德·桑蒂拉纳（Giorgio de Santillana）与法兰克福大学教授赫尔塔·冯·德亨德（Hertha von Dechend）提出，古代神话并非随机的虚构故事，而是对天体宇宙学（celestial cosmology）的准确描述，以寓言形式传承下来以便口述传播。他们认为，神话中的混沌、怪兽与暴力象征着宇宙各个时代的形成力量。在其著作《哈姆雷特的磨坊：神话与时间结构的论文》（Hamlet’s Mill: An Essay on Myth and the Frame of Time, 1969）中，两人主张古代神话并无历史或事实基础，而是以宇宙学形式编码的天文现象，尤其是岁差（precession of the equinoxes）。\(^\text{[147]}\)不过，这一理论在学术界并未得到广泛接受。
\subsubsection{权力的展现}
\begin{figure}[ht]
\centering
\includegraphics[width=6cm]{./figures/2702dbdfaa401567.png}
\caption{} \label{fig_ArmY_12}
\end{figure}
通过在服饰中融入天体图案，佩戴者便可宣称：他们在尘世间的权力，源自于天上的力量。正如学者所指出的，荷马在《伊利亚特》中所描述的阿喀琉斯之盾（Shield of Achilles），实际上也是一份星座目录\(^\text{[148]}\)。在北美地区，科曼奇族（Comanche）的岩刻图像中所描绘的盾牌，似乎也包含了与金星（Venus）相关的象征\(^\text{[149]}\)。

至日（solstice）方向的建筑对齐，同样可以被视为权力的展示。在的的喀喀湖（Lake Titicaca）上的太阳岛（Island of the Sun）——传说中太阳诞生的地方——从岛上的祭祀广场望去，六月至日（冬至）清晨的太阳会从附近山脊上的两座塔之间升起。岛上神圣的区域被石墙与其余部分隔离开来。民族志记录表明，只有印加统治阶层（Inca ruling elite）才被允许进入这片神圣的空间。普通朝圣者则必须站在仪式区外的平台上，从远处观看太阳升起、光线穿越双塔的景象\(^\text{[150]}\)。

在埃及（Egypt），卡纳克的阿蒙神庙（Temple of Amun-Re at Karnak）长期以来都是考古天文学（archaeoastronomy）研究的重点。研究表明，若考虑到黄道倾角（obliquity of the ecliptic）随时间的变化，这座大神庙的中轴线与冬至日太阳升起的方向严格对齐\(^\text{[151]}\)。神庙长廊的长度经过精确设计，使得只有在特定季节的清晨阳光才可照入神庙深处，其余时间则保持黑暗。

在稍后的时期，亚历山大的塞拉比昂神殿（Serapeum of Alexandria）据说也拥有类似的太阳光轴结构：在特定的日出时刻，一束阳光会穿过殿堂，正好掠过塞拉比斯（Serapis）神像的唇部，从而象征太阳亲吻神灵的神圣时刻\(^\text{[152]}\)。
\subsection{全球性的天文遗址}
考古天文学家克莱夫·拉格尔斯（Clive Ruggles）与米歇尔·科特（Michel Cotte）最近合编了一部关于天文与考古天文学遗产地（heritage sites of astronomy and archaeoastronomy）的著作，书中讨论了全球范围内具有代表性的天文遗址，并提出了考古天文学遗址分类的标准与方法\(^\text{[153]}\)。

正如天文学家埃德温·克鲁普（Edwin C. Krupp）所言：

在英格兰的巨石阵（Stonehenge）、法国的卡纳克（Carnac）、埃及与尤卡坦（Yucatán）——乃至全球各地——都存在那些充满神秘意义的古代遗迹。这些纪念碑具有深刻的天文学意义。它们所体现的，是与我们登月、以及将飞船送上火星时所展现的——同样的探索精神与信念。\(^\text{[154]}\)
\subsubsection{纽格兰奇}
\begin{figure}[ht]
\centering
\includegraphics[width=6cm]{./figures/e4dca096df892f6d.png}
\caption{} \label{fig_ArmY_13}
\end{figure}
纽格兰奇（Newgrange）是一座位于爱尔兰共和国境内的甬道墓（passage tomb），建造年代约为公元前 3300 至公元前 2900 年之间\(^\text{[155]}\)。在冬至（Winter Solstice）前后数日的清晨，阳光会沿着墓道的中央通道直射进墓室深处的核心位置。

这一现象的独特之处并不在于光线能够进入墓道，而在于它并非通过主入口射入。相反，阳光是通过主门上方的一个中空结构——即所谓的“天窗盒”（roof-box）进入的，这一结构是由考古学家迈克尔·奥凯利（Michael O’Kelly）在发掘过程中发现的\(^\text{[156]}\)。正是这一“天窗盒”的存在，强烈地表明该墓在建造时被赋予了明确的天文取向（astronomical aspect）。

在关于考古天文学遗址可信度的讨论中，米歇尔·科特（Michel Cotte）与克莱夫·拉格尔斯（Clive Ruggles）将纽格兰奇列为“普遍接受”（Generally accepted）级别的遗址——这是他们所定义的四个可信度等级中最高的一级\(^\text{[81]}\)。

克莱夫·拉格尔斯（Clive Ruggles）指出：

……几乎没有人——无论是考古学家还是天文学家——会怀疑，这座纪念碑在设计之初就被有意地赋予了强烈的天文学象征意义。这一事实清楚地表明，天文学与葬仪仪式之间的联系，至少值得进一步深入探究。\(^\text{[117]}\)
\subsubsection{埃及}
\begin{figure}[ht]
\centering
\includegraphics[width=6cm]{./figures/e2ad48c54af60b82.png}
\caption{吉萨金字塔群} \label{fig_ArmY_15}
\end{figure}

\begin{figure}[ht]
\centering
\includegraphics[width=6cm]{./figures/24bac962fd900cb5.png}
\caption{猎户座腰带叠加在吉萨金字塔群上的示意图，用以说明“猎户对应理论”（Orion Correlation Theory）。从左至右依次为：大金字塔（胡夫金字塔）对应恒星参宿一（Alnitak），卡夫拉金字塔对应恒星参宿二（Alnilam），孟卡拉金字塔对应恒星参宿三（Mintaka）。} \label{fig_ArmY_16}
\end{figure}
自弗林德斯·皮特里（Flinders Petrie）首次以现代方法精确测量吉萨金字塔群的方位以来，学者们提出了多种关于这些金字塔原始定向方式的天文学假说\(^\text{[157][158][159]}\)。最新的研究提出，这些定向可能是通过观测“大熊座”（Plough / Big Dipper）中两颗恒星的位置来实现的——在古埃及文化中，该星座被称为“牛腿”（the thigh）。研究者推测，古埃及人可能利用这两颗恒星的垂直连线来确定北方：当两星在天空中垂直排列时，通过垂球（plumb bob）进行目测校准，即可确定真北方向。基于这种模型计算出的方位偏差，与目前普遍接受的各金字塔建造年代在时间上高度吻合\(^\text{[160]}\)。这一发现为古埃及工程师如何在无现代仪器的情况下实现如此精确的定向提供了令人信服的天文学解释，也揭示了古代文明对恒星运动规律的深刻理解。
\begin{figure}[ht]
\centering
\includegraphics[width=6cm]{./figures/cd0abdec3ed640b5.png}
\caption{塞涅姆特墓室天文穹顶上的星座图} \label{fig_ArmY_14}
\end{figure}
一些学者\(^\text{[161][162]}\)认为金字塔的布局对应猎户座腰带上的三颗恒星，但这一理论已受到权威天文学家的批评。实际上，该遗址的布局更可能受一种壮观的神圣显现（hierophany）所支配——即在夏至这一天，从狮身人面像（Sphinx）平台上观察，太阳与两座巨大的金字塔共同形成象征“地平线”的符号“Akhet”，这也是大金字塔（Great Pyramid）的名称。此外，三座金字塔的东南角均朝向赫利奥波利斯神庙（Temple of Heliopolis），此现象最早由埃及学家马克·莱纳（Mark Lehner）发现。

约公元前1470年的塞涅姆特（Senenmut）墓室天文穹顶上绘有“天体图”（Celestial Diagram），其中以圆盘形式描绘了绕极星座。每个圆盘被分为24个部分，暗示出24小时的时间划分。星座被描绘为埃及的神祇，同时也明显展示了对月相周期的观测。
\subsubsection{埃尔卡斯蒂略神庙}
埃尔卡斯蒂略（El Castillo），又称“库库尔坎金字塔”（Kukulcán's Pyramid），是一座位于墨西哥玛雅遗址奇琴伊察中心的中美洲阶梯式金字塔。其建筑特征显示出显著的天文学元素。金字塔四面各有一条通往顶端的阶梯，每条共有91级台阶；再加上最上方平台的一阶，总计为365级。这一数字可能象征着一年中的365天（365.25），或代表在10,000次地球自转中月球的轨道数（365.01）。
